\documentclass{article}
\usepackage{import}
\import{../template/}{article}
\graphicspath{{../../images}}
\addbibresource{../template/library.bib}

\title{Manifold Theory}
\author{PENG}
\date{Created: 20260324, Version: 20260627}

\begin{document}
\maketitle
\tableofcontents
\section{Set Theory}
\subsection{Relations}
\begin{definition}
    Equivalence relation $\sim$:

    \begin{itemize}
        \item reflexive: $x \sim x$
        \item symmetric: $x \sim y \Rightarrow y \sim x$
        \item transitive: $x \sim y, y \sim z \Rightarrow x \sim z$
    \end{itemize}
\end{definition}

\begin{definition}
    Equivalence class of each $x \in X$:

    $[x] = \{y \in X : y \sim x\}$
\end{definition}

\begin{definition}
    The collection of all $[x]$ is $X/\sim$
\end{definition}

\begin{example}
    Let $A = \{1,2,3\}$, $R = \{(x,y) : x = y\}$ on $A$

    \begin{itemize}
        \item reflexive: $x=x$
        \item symmetric: $x=y \Rightarrow y=x$
        \item transitive: $x=y,y=z \Rightarrow x=z$
    \end{itemize}

    Thus $R$ is an equivalence relation.

    Equivalence class:

    $[1] = \{1\}$

    $[2] = \{2\}$

    $[3] = \{3\}$
\end{example}

\begin{example}
    Let $A = \{1,2,3\}$, $R = \{(x,y) : x + y \text{ is even}\}$ on $A$

    \begin{itemize}
        \item reflexive: $x+x$ is even
        \item symmetric: $x+y \text{ is even} \Rightarrow y+x \text{ is even}$
        \item transitive: $x+y \text{ is even}, y+z \text{ is even} \Rightarrow \begin{cases}
                      x \text{ is even}, y \text{ is even}, z \text{ is even} \\
                      x \text{ is odd}, y \text{ is odd}, z \text{ is odd}
                  \end{cases} \Rightarrow x+z \text{ is even}$
    \end{itemize}

    Thus $R$ is an equivalence relation.

    Equivalence class:

    $[1] = \{1,3\}$

    $[2] = \{2\}$

    $[3] = \{1,3\}$
\end{example}

\begin{definition}
    A partition of $X$ is a collection of disjoint nonempty subsets of $X$ whose union is $X$.
\end{definition}

\begin{example}
    $X = \{1,2,3\}$

    $P = \{\{2\},\{1,3\}\}$

    $P = \{\{1\},\{3\},\{2\}\}$

    $P = \{\{1,2,3\}\}$

    $\{\{1\},\{2\},\{\}\}$: nonempty

    $\{\{1\},\{2\}\}$: union is not $X$

    $\{\{1,3\},\{2,3\}\}$: disjoint

    $\{\{1,3\},\{2,4\}\}$: union is not $X$
\end{example}

\begin{exercise}
    Collection of equivalence class $\Rightarrow$ partition

    Partition $\Rightarrow$ unique equivalence relation
\end{exercise}

\begin{proof}
    $X/\sim \Rightarrow$ partition:

    For $\forall x_{1}, x_{2} \in X:$

    $[x_{1}]=\{y \in X: y \sim x_{1}\}$

    $[x_{2}]=\{y \in X: y \sim x_{2}\}$


    If $x_{1} \nsim x_{2}, \forall y \in [x_{1}] \Rightarrow y \notin [x_{2}] \Rightarrow$ disjoint nonempty subsets

    If $x_{1} \sim x_{2}, \forall y \in [x_{1}] \Rightarrow y \in [x_{2}] \Rightarrow [x_{1}] = [x_{2}]$
\end{proof}

\begin{proof}
    Partition $\Rightarrow \sim$
\end{proof}

\section{Euclidean Spaces}

\begin{equation*}
    x =
    \begin{bmatrix}
        x_{1}  \\
        x_{2}  \\
        \vdots \\
        x_{n}
    \end{bmatrix}
    \in \mathbb{R}^{n}
\end{equation*}

\begin{definition}
    For vectors:

    \begin{equation*}
        x \cdot y = x^{T}y = x_{1}y_{1} + x_{2}y_{2} + \dots + x_{n}y_{n}
    \end{equation*}
\end{definition}

\begin{definition}
    Norm of $x$:

    \begin{equation*}
        |x| = \sqrt{(x_{1})^2 + (x_{2})^2 + \dots + (x_{n})^2}
    \end{equation*}
\end{definition}

\begin{exercise}
    \begin{equation*}
        \max \{|x_{1}|,|x_{2}|,\dots,|x_{n}|\} \leq |x| \leq \sqrt{n} \max \{|x_{1}|,|x_{2}|,\dots,|x_{n}|\}
    \end{equation*}
\end{exercise}

\begin{proof}
    Let $x_{m} = \max \{|x_{1}|,|x_{2}|,\dots,|x_{n}|\}$:

    \begin{equation*}
        \sqrt{(x_{m})^2} \leq \sqrt{(x_{1})^2 + \dots + (x_{m})^2 + \dots + (x_{n})^2} \leq \sqrt{(x_{m})^2 + \dots + (x_{m})^2}
    \end{equation*}
\end{proof}

\begin{definition}
    $\theta$ between $x$ and $y$:
\end{definition}

\section{Topological Spaces}
Topology \autocite{LEE_2010_INTRODUCTION_TO_TOPOLOGICAL_MANIFOLDS}

Four mechanisms (normal, segment, merge, torus) exist in every fluid point: fractal.

\begin{figure}[h]
    \centering
    \includegraphics{sloshing-topology.png}
    \caption{Sloshing topology.}
    \label{fig:sloshing-topology}
\end{figure}

\printbibliography
\end{document}
